\documentclass{article}
\newif\ifdraft
\drafttrue
\usepackage{graphicx} % Required for inserting images
\usepackage{hyperref}
\usepackage{booktabs}
\newcommand{\manual}[1]{\texttt{#1}}





\title{Laboratory Exercise 1: Dynamical Pendulum}
\author{Magnus Støren Wedèn, Jon Magnus Vårheim}
\date{September 2026}

\begin{document}

\maketitle

\section{Introduction}
\subsection{theory}


\section{method}


\section{First watch observations:}
\textit{Her har vi gjort veldig macro-vis observasjon av videoen. Denne delen skal inneholde forklaring av hva vi ser på.
Eks:
SP1 kan observeres i dette tidsrommet. Her ser vi pendulumen oscillerer om dette punktet. Kaya dytter på pendulumen. etc etc
SP2 kan observes i dette tidsrommet. Her er en ytre kraft eller kontroller som påvirker systemet slik at pendulumen balanaseres om sitt nye setpoint. etc etc}

- Settle point 1. 8s - 38.5s\\
- Settle point 2. 53s - 1:09 (controlled upright balancing)\\
- Disturbances 1:54 - 2:01\\

- Disturbance events time stamps\\
1:55.215-248\\
1:56.083-116\\
1:56.850-950\\
1:57.784-884\\
2:00.354-454\\

Paste this in place of the existing **“Video and available data”** subsection and interval table. It uses the packages and `\manual{}` command already in your template.

**Assumption:** `1:55.215–248` means `00:01:55.215–00:01:55.248`. The five distinct event windows are retained below.

The wording uses **Stable Point** and separates visible behaviour from possible explanations. Your approximate SP1 endpoint, **38.5 s**, differs from the exported boundary, **38.885 s**; the table uses the export. It also includes your second SP2 interval.

```latex
\subsection{Video overview and selected analysis intervals}
\label{sec:video-overview}

An initial qualitative review of the video was used to identify
the main phases of the experiment before taking detailed
measurements. The review focused on the overall motion of the
pendulum and cart, changes in operating condition, and possible
disturbance events. This follows the observation procedure
described in the laboratory manual. \manual{3--4}

The recording was divided into intervals covering natural motion
around Stable Point 1 (SP1), controlled operation around Stable
Point 2 (SP2), and disturbances during upright operation.
SP1 denotes the naturally stable downward equilibrium, whereas
SP2 denotes the upright equilibrium maintained by feedback
control. \manual{2, 6--8}

For the SP1 intervals, the analysis considers displacement and
release of the pendulum, its subsequent oscillation about the
downward equilibrium, and any visible cart motion.
For the SP2 intervals, the analysis considers angular deviations
from upright and the associated corrective cart movements.
The disturbance interval is examined more closely to identify
changes in motion, the sequence of visible responses, and the
subsequent recovery or loss of stabilisation.
\manual{10--11, 14--20}

Table~\ref{tab:intervals} lists the selected analysis intervals.
Timestamps refer to the original recording and use the format
hh:mm:ss.sss. The timestamps and frame boundaries were obtained
from the segment exports and matched by segment name.
The interval descriptions identify the intended focus of the
analysis; they do not establish the cause of an observed event.

\begin{table}[htbp]
    \centering
    \small
    \caption{Selected analysis intervals in the original video,
    with timestamps and exported frame boundaries.}
    \label{tab:intervals}
    \begin{tabular}{@{}lccrr@{}}
        \toprule
        Analysis interval
        & Start time
        & End time
        & \shortstack{Start\\frame}
        & \shortstack{End\\frame} \\
        \midrule
        SP1: initial observation
        & 00:00:08.041 & 00:00:11.812 & 241 & 354 \\

        SP1: input/disturbance
        & 00:00:12.980 & 00:00:15.452 & 389 & 463 \\

        SP1: input/disturbance continued
        & 00:00:20.287 & 00:00:23.790 & 608 & 713 \\

        SP1: oscillation
        & 00:00:33.188 & 00:00:38.885 & 995 & 1165 \\

        SP2: controlled operation 1
        & 00:00:53.020 & 00:01:09.937 & 1589 & 2096 \\

        SP2: controlled operation 2
        & 00:01:13.974 & 00:01:22.616 & 2217 & 2476 \\

        SP2: disturbance sequence
        & 00:01:54.015 & 00:02:01.197 & 3417 & 3632 \\
        \bottomrule
    \end{tabular}

    \smallskip
    \begin{minipage}{\linewidth}
        \footnotesize
        \textit{Source:} Timestamp and frame segment exports
        for the laboratory video.
        \textit{Note:} The exports do not specify whether
        frame numbering begins at zero or one, or whether
        the end frame is included. Frame values are therefore
        reported as exported boundary indices.
    \end{minipage}
\end{table}

\subsubsection{Candidate disturbance-event windows}
\label{sec:disturbance-windows}

Within the selected disturbance sequence, five candidate
event windows were marked for closer inspection, as listed
in Table~\ref{tab:disturbance-windows}.
These windows guide the frame-by-frame analysis.
They are preliminary event locations, rather than confirmed
disturbance durations or onset-time uncertainty intervals.

For each candidate event, direct observations are recorded
before a disturbance mechanism is proposed. The analysis
considers which component responds first, whether corrective
cart motion remains visible, and whether the pendulum returns
to upright operation. Where the evidence does not support a
reliable explanation, the disturbance is classified as unknown.
\manual{3--4, 17--20}

\begin{table}[htbp]
    \centering
    \small
    \caption{Preliminary disturbance-event windows in the
    original video. Timestamps use hh:mm:ss.sss.}
    \label{tab:disturbance-windows}
    \begin{tabular}{@{}ccc@{}}
        \toprule
        Candidate event & Window start & Window end \\
        \midrule
        D1 & 00:01:55.215 & 00:01:55.248 \\
        D2 & 00:01:56.083 & 00:01:56.116 \\
        D3 & 00:01:56.850 & 00:01:56.950 \\
        D4 & 00:01:57.784 & 00:01:57.884 \\
        D5 & 00:02:00.354 & 00:02:00.454 \\
        \bottomrule
    \end{tabular}
\end{table}
```


\section{Second watch observations:}



\section{results}
noen notater:
Section 1:
x(t) = 0 - cart i midtpunkt \newline
x´(t) = 0 - cart står stille \newline
x´´(t) = 0 - cart aksellerer ikke \newline
\(\phi\)(t) = 0 - pendulum i SP1, ingen radiell displacement \newline
\(\phi\)´(t) = 0 - Pendulum i SP1, stasjonær, ingen velocity \newline
\(\phi\)´´(t) = 0 - pendulum har ingen radiell aksellerasjon \newline
M = ? \SI{1.12}{kg} - masse til cart \newline
m = ?\SI{0.120}{kg} - masse til pendulum \newline
l = ? \SI{ 0.0139231}{} \newline
Jp = \newline
bp = \newline
g = 9.8m/\(s^2\) \newline
\(\phi\)ref(t) = SP1 = \(0\circ \) = 0 rad \newline
Xref(t) == x(t) = 0 - cart i midtpunkt\newline





\section{analysis}

\section{discussion}
uncertainty of starting point

-angle/wiev
-fps
-cart position


\textbf{Part 1: SP1 - Questions and Answers}

\begin{enumerate}

 \item What visual evidence indicates that SP1 is a stable equilibrium?

 That the pendulum returns toward the downward position after being displaced. No external control is needed to bring it back toward equilibrium.

 \item Why does the pendulum pass through SP1 rather than stopping immediately?

 Because the pendulum has kinetic energy. This makes it continue past the equilibrium position before reversing direction.

 \item Classify the response as underdamped, critically damped or overdamped. Support the
classification using the observed motion

The response is underdamped: the pendulum oscillates several times before coming to rest. The amplitude gradually decreases with time.

\item Why does the oscillation amplitude decrease with time?

The mechanical energy is lost due to friction and air resistance.

\item Which physical effects are likely to dissipate mechanical energy?

Mechanical energy is dissipated through pivot friction, rolling friction in the cart, bearing friction, and air resistance.

\item Does the oscillation period remain approximately constant? Support the answer using
the calculated periods.

Yes. The calculated periods are similar, and the small differences are within the measurement uncertainty. Therefore, the period can be considered approximately constant.

\item Are the successive amplitude ratios approximately constant?

Yes. The calculated amplitude ratios are close to each other, this means that the oscillation amplitude decreases at a fairly constant rate.

\item Are the positive and negative peak displacements approximately symmetric? If not,
suggest possible explanations.

The peaks are not perfectly symmetric. This may be caused by friction, slight measurement errors, camera perspective, or small movements of the cart.

\item Does any cart motion affect the pendulum response and how?

Yes. As the pendulum swings, it makes the cart move slightly. This gives energy from the pendulum to the cart and increases the damping.

\item How do the video frame rate, camera perspective, encoder resolution and choice of
equilibrium reference affect the measurement uncertainty?

A low frame rate makes it harder to determine exact timestamps. Camera perspective can make angles appear different. Image quality also affect angle measurements. 

\item Does the pendulum return to precisely the same final position from which it was
released? Explain any observed difference

No. Small differences can occur because of friction, measurement uncertainty, and mechanical imperfections in the system.

\item Which observed-motion aspects can be omitted by a simple ideal pendulum model?

An ideal pendulum model ignores air resistance, friction, cart motion, measurement noise, and other real-world effects. It assumes a fixed pivot and no energy losses.

 \end{enumerate}

 \textbf{Part 2: SP2 - Questions and Answers}

\begin{enumerate}

 \item What evidence indicates that the feedback controller is active?

 \item

\end{enumerate}



\section{conclusion}

\section{}
www.when2meet.com:

\url{https://www.when2meet.com/?38389767-u2wP9}
\section{references}
\section{attachments}
\end{document}
